\documentclass[11pt]{article}
\usepackage[T1]{fontenc}
\usepackage{lmodern}
\usepackage[margin=1in]{geometry}
\usepackage{amsmath,amssymb,amsthm,mathtools}
\usepackage{booktabs,longtable,array}
\usepackage{microtype}
\usepackage{enumitem}
\usepackage{xurl}
\usepackage[hidelinks]{hyperref}
\hypersetup{
 pdftitle={Central Poupard Numbers: OEIS Identities and Shifted-Moment Transitions},
 pdfauthor={Research article prepared for Vladimir Reshetnikov},
 pdfsubject={Continued fractions, asymptotic sectors, coexistence, complex zeros, and inversion},
 pdfkeywords={A125054, A000182, Poupard numbers, tangent numbers, asymptotics, Lambert W}
}
\usepackage{fancyhdr}
\usepackage{titlesec}
\usepackage{verbatim}
\setlength{\headheight}{14pt}
\pagestyle{fancy}
\fancyhf{}
\fancyhead[L]{\small Central Poupard numbers and shifted moments}
\fancyhead[R]{\small Research article}
\fancyfoot[C]{\thepage}
\renewcommand{\headrulewidth}{0.3pt}
\setlength{\parindent}{1.2em}
\setlength{\parskip}{0.25em}
\setlength{\emergencystretch}{2em}
\allowdisplaybreaks[2]
\titleformat{\section}{\large\bfseries}{\thesection}{0.7em}{}
\titleformat{\subsection}{\normalsize\bfseries}{\thesubsection}{0.7em}{}
\newtheorem{theorem}{Theorem}[section]
\newtheorem{proposition}[theorem]{Proposition}
\newtheorem{lemma}[theorem]{Lemma}
\newtheorem{corollary}[theorem]{Corollary}
\theoremstyle{definition}
\newtheorem{definition}[theorem]{Definition}
\newtheorem{question}{Research question}
\theoremstyle{remark}
\newtheorem{remark}[theorem]{Remark}
\newcommand{\A}{\mathcal A}
\newcommand{\E}{\mathbb E}
\newcommand{\Prob}{\mathbb P}
\newcommand{\R}{\mathbb R}
\newcommand{\Z}{\mathbb Z}
\newcommand{\dd}{\,\mathrm d}
\newcommand{\sech}{\operatorname{sech}}
\newcommand{\fall}[2]{(#1)_{\underline{#2}}}
\newcommand{\Sfrac}{\mathsf S}
\newcommand{\Jfrac}{\mathsf J}
\newcommand{\coef}[2]{[#1^{#2}]}
\newcommand{\N}{\mathbb N}
\newcommand{\Normal}{\mathcal N}
\newcommand{\eps}{\varepsilon}
\newcommand{\st}{\star}
\newcommand{\ind}{\mathbf 1}
\makeatletter
\newcommand{\tableinput}[1]{\@@input #1 }
\makeatother
\newcommand{\smallurl}[1]{{\small\url{#1}}}
% Collection edition (batch 73O1): notes added when the report was written
% into ProveIt's research-report collection.
\newcommand{\writenote}[1]{\par\smallskip\noindent
 \textit{[Write note, batch 73O1, 1 October 2026.]}\ #1\par\smallskip}
\title{\textbf{Central Poupard Numbers}\\[0.35em]
\Large Two OEIS conjectural statements, complete asymptotic sectors,\\
and a universal transition for shifted moments}
\author{Research article prepared for Vladimir Reshetnikov}
\date{1 October 2026}

\begin{document}
\maketitle
\begin{abstract}
Let $T_n=(2n+1)![z^{2n+1}]\tan z$ and let
$A_n(c)=\sum_{k=0}^n\binom nk c^{n-k}T_k$.
The specialization $A_n(1)$ is OEIS A125054, the central Poupard numbers.
We prove the continued-fraction identity recorded there as a conjecture in
December 2025, and prove the stronger congruence $A_n(1)\equiv3\pmod9$
for every $n\ge1$, which settles the entry's valuation observation.
Both proofs use elementary formal continued fractions.

The main analytic development is considerably broader.  A positive integral
representation yields an exact, convergent decomposition into odd-pole
exponential sectors, an explicit remainder bound, and inverse-power expansions
of every fixed sector to arbitrary order.  We then prove a uniform, all-order
endpoint--saddle expansion when the shift is $c=\tau n^2$.  Its first-order
transition occurs at $\tau_*=0.262466543368419\ldots$, but its finite-size
coexistence window is displaced by a multiple of $(\log n)/n$.
We obtain the limiting mixture law and the Gaussian fluctuations inside its
macroscopic component.  Positivity also yields a complex coexistence limit
and a rigorously located array of polynomial zeros.  A general theorem
extends the transition to moments
of $y^p$ under any probability density with tail $Cy^\beta e^{-\alpha y}$,
$p>1$.  Finally, we invert the fixed-shift growth using Lambert $W$, including
an exponentially small displacement relative to an exact one-sector inverse.

All essential arguments are supplied.  Classical generating functions,
continued-fraction machinery, and the previously recorded leading asymptotic
are explicitly distinguished from the extensions developed here.
The accompanying code provides exact finite checks and high-precision
numerical diagnostics; it is not a Lean formalization or an interval proof.
Worldwide publication priority for the broader extensions has not been
established.
\end{abstract}

\writenote{This is the collection edition of the report, written into
ProveIt's research-report collection from manuscript~47 of batch~73O1
(provenance in Appendix~\ref{pou:app:edition}). Every label carries the
prefix \texttt{pou:}; the \texttt{Makefile} now lives in \texttt{code/},
and \texttt{requirements.txt} and \texttt{BUILD\_REPORT.json} in
\texttt{data/}. Text added in the write is marked like this note;
everything else is the delivered text. The OEIS statements quoted here
(Bala's conjectured fraction, Chapoton's observation, Kotesovec's
leading term) are as the authors inspected the entry on 1~October 2026;
the intake did not re-fetch it.}

\clearpage
\begingroup
\small
\setlength{\parskip}{0pt}
\tableofcontents
\endgroup
\clearpage

\section{Scope, provenance, and the central results}\label{pou:sec:scope}

The two concrete OEIS targets are the continued fraction labelled
``Conjectured g.f.'' by Peter Bala on 15 December 2025, and the observation
by F. Chapoton on 2 August 2021 that the noninitial terms of A125054 appear
to have $3$-valuation one.  Both labels were present in the entry inspected
for this article \cite{oeis125054}.  The statements proved below are stronger
than finite verification: the continued fraction is an identity of formal
power series, and the congruence holds for every positive index.
These are resolutions of the displayed statements, not a claim that no
unpublished or unlocated proof previously existed.

The sequence is not introduced merely to manufacture an asymptotic example.
It is the center of the Poupard triangle, part of the finite-difference
calculus of tangent and secant enumeration \cite{foatahan,oeis125053}.
Its tangent-number representation makes it a particularly explicit test
case for the competition between a large deterministic shift and the
factorial growth of moments.

The present work continues the repository's emphasis on exact identities,
all-order expansions, inversion, and transparent verification boundaries.
The inspected ProveIt snapshot is commit
\begin{center}
\texttt{ef66df0fdb8fc64cbc80756d622cab0e8e85e262}.
\end{center}
In particular, the repository's logarithmic-endpoint/Lambert-chart report
\cite{repochart} is a methodological precedent, not a theorem used as a
hypothesis here.  Its countable-action model is different from the shifted
moment integrals below.  No repository result, formalized or otherwise,
is required for the proofs in this article.  The repository was not modified.

\writenote{The precedent cited as \cite{repochart} is the transseries
package
\path{Analysis/Transseries/docs/series-and-transseries/Logarithmic_Critical_Endpoint_Lambert_Charts/};
nothing is imported from it. At the repository's current state no other
report or formal development treats A125054, A125053, A000182 or the
shifted tangent moments $A_n(c)$. One constant recurs: the transition
parameter $s_*=2+W_0(-2e^{-2})$ of \eqref{pou:eq:starintro} is the
number called $\tau$ in the collection report on A277364
\cite{bellreport}, where it solves $\tau/(1-e^{-\tau})=2$ (the same
equation $2-s=2e^{-s}$) and governs the omitted tail $B_n-a(n)$ of the
truncated Stirling sum $a(n)=\sum_{k\le n/2}S(n,k)$. The objects differ;
see the note after the numerical values of $s_*$ and $\tau_*$ in
Section~\ref{pou:sec:geometry}.}

\subsection{What is known, what is proved here}

The tangent generating function, its classical Stieltjes fraction, the
Poupard bivariate generating function, and the leading equivalent of
A125054 are prior results \cite{bala,foatahan,oeis125054}.  The moment and
Hankel consequences are included because they explain the structure; they
are not advertised as discoveries.  Our substantive extension consists of
the exact-sector remainder control, the uniform two-contribution expansion,
the moving coexistence window and its probability laws, the $p>1$
universality theorem, the associated complex zero array, and the inversion
formulas.
These are proved in the article, but a targeted source search is not an
exhaustive priority determination.

For orientation, the central statements are as follows.  Put
\begin{equation}\label{pou:eq:definition}
 \alpha=\frac\pi2,\qquad
 T_n=(2n+1)!\,[z^{2n+1}]\tan z,\qquad
 A_n(c)=\sum_{k=0}^n\binom nk c^{n-k}T_k,\qquad a_n=A_n(1).
\end{equation}
The structural proofs establish the conjectured fraction and
\begin{equation}\label{pou:eq:congruenceintro}
 a_n\equiv3\pmod9\quad(n\ge1).
\end{equation}
For fixed $c\ge0$ we obtain the exact sector expansion
\begin{equation}\label{pou:eq:sectorintro}
 A_n(c)=\frac{2\Gamma(2n+2)}{\alpha^{2n+2}}
 \sum_{\substack{q\ge1\\q\ \mathrm{odd}}}
 q^{-2n-2}H_n(c\alpha^2q^2),\qquad
 H_n(z)=\sum_{j=0}^n\frac{(z/4)^j}{j!\fall{n+\tfrac12}{j}}.
\end{equation}
For the large shift $c=\tau n^2$, the transition constant is characterized by
\begin{equation}\label{pou:eq:starintro}
 2-s_*=2e^{-s_*},\quad 1<s_*<2,\qquad
 \tau_* =\frac{s_*(2-s_*)}{\alpha^2}.
\end{equation}
In a precisely specified moving window, the normalized moment tends to
$1+e^{-w}$, rather than to a single saddle contribution.
The coefficient procedures and uniform errors make these formulas usable
beyond their leading terms.

\subsection{Conventions about asymptotics and formal series}

The symbol $\sim\sum_{r\ge0}c_r n^{-r}$ means a Poincar\'e expansion:
for every fixed $M$, truncation before $n^{-M}$ has error $O(n^{-M})$.
It does not assert convergence of the infinite inverse-power series.
An exact positive sector decomposition and an inverse-power expansion of
one sector are different assertions and will be kept separate.

All ordinary generating functions of the factorially growing sequences in
this article are formal power series.  Their radius of convergence is zero.
An infinite continued fraction means its coefficientwise stabilized formal
expansion, not its evaluation at a small positive real number.  All
probability statements concern explicitly normalized positive weights.

\section{From the Poupard triangle to tangent moments}\label{pou:sec:poupard}

The first rows of the relevant Poupard triangle are
\[
\begin{array}{c}
1,\\
1,\ 3,\ 1,\\
5,\ 15,\ 21,\ 15,\ 5,\\
61,\ 183,\ 285,\ 327,\ 285,\ 183,\ 61.
\end{array}
\]
To fix indexing without relying on these examples, write $h_{i,j}$ for its
entry with $i+j=2n$, $i,j\ge0$, and set $h_{i,j}=0$ when $i+j$ is odd.
The classical bivariate exponential generating function is
\begin{equation}\label{pou:eq:rhombus}
 Z(x,y)=\sum_{i,j\ge0}h_{i,j}\frac{x^iy^j}{i!j!}
       =\cos(x-y)\sec^2(x+y).
\end{equation}
This is the secant case of Foata and Han's generating-function theorem
\cite[Theorem 1.2]{foatahan}.

For completeness, the algebraic content of this representation can also be
checked directly.  It is symmetric in $x,y$, its boundary value is
$Z(x,0)=\sec x$, and it satisfies
\begin{equation}\label{pou:eq:pde}
 (\partial_x-\partial_y)^2Z=-4Z.
\end{equation}
Consequently its coefficients satisfy
\[
 h_{i+2,j}-2h_{i+1,j+1}+h_{i,j+2}=-4h_{i,j}.
\]
The two boundary values of each even row and its prescribed second
differences determine that row uniquely: the difference of two solutions
would be an affine function of the position, vanishing at both ends.
Thus \eqref{pou:eq:rhombus} characterizes the triangle with secant-number edges.
The combinatorial interpretation is prior work; the following short
operator calculation explains the central specialization used here.

\begin{proposition}\label{pou:prop:central}
The central entry is $h_{n,n}=A_n(1)$.  More generally, the diagonal
exponential coefficient of
\[
 Z_c(x,y)=\cos\bigl(\sqrt c\,(x-y)\bigr)\sec^2(x+y)
\]
is $A_n(c)$; the expression is a formal power series in $c$, so the square
root creates no ambiguity.
\end{proposition}
\begin{proof}
Set $s=x+y$ and $d=x-y$.  Then
$\partial_x\partial_y=\partial_s^2-\partial_d^2$.
Since $\partial_d^2\cos(\sqrt c\,d)=-c\cos(\sqrt c\,d)$,
\[
 (\partial_x\partial_y)^nZ_c(0,0)
 =\left.(\partial_s^2+c)^n\sec^2s\right|_{s=0}
 =\sum_{k=0}^n\binom nk c^{n-k}T_k.
\]
Here $(\sec^2s)^{(2k)}(0)=(\tan s)^{(2k+1)}(0)=T_k$.
The left side is exactly $n!^2[x^ny^n]Z_c(x,y)$.
\end{proof}

One exact recurrence useful for checking the normalization comes from
$\tan'=1+\tan^2$:
\begin{equation}\label{pou:eq:tangentrec}
 T_0=1,\qquad
 T_n=\sum_{j=0}^{n-1}\binom{2n}{2j+1}T_jT_{n-1-j}\quad(n\ge1).
\end{equation}
It gives $T_n=1,2,16,272,7936,\ldots$ (OEIS A000182~\cite{oeis182}), hence
$a_n=1,3,21,327,9129,\ldots$.

\section{A proof of the conjectured continued fraction}\label{pou:sec:cf}

\subsection{Formal contractions and binomial translation}

For a sequence of weights $\gamma_1,\gamma_2,\ldots$, write
\[
 \Sfrac(\gamma;x)=
 \cfrac1{1-\cfrac{\gamma_1x}{1-\cfrac{\gamma_2x}{1-\cdots}}}.
\]
For diagonal weights $b_0,b_1,\ldots$ and downward weights
$\lambda_1,\lambda_2,\ldots$, write
\[
 \Jfrac(b,\lambda;x)=
 \cfrac1{1-b_0x-\cfrac{\lambda_1x^2}
 {1-b_1x-\cfrac{\lambda_2x^2}{1-b_2x-\cdots}}}.
\]
These definitions make sense over any commutative coefficient ring.
At any given degree only finitely many levels can contribute.
The elementary contraction and translation rules are classical
\cite{bala}; we record their exact indexing.

\begin{lemma}\label{pou:lem:contractions}
The following identities hold formally.
\begin{enumerate}[label=(\roman*),itemsep=2pt]
\item The even contraction of $\Sfrac(\gamma;x)$ has
\[
 b_0=\gamma_1,\qquad b_k=\gamma_{2k}+\gamma_{2k+1}\ (k\ge1),
 \qquad \lambda_k=\gamma_{2k-1}\gamma_{2k}.
\]
\item If $S=\Sfrac(\gamma;x)$, then $(S-1)/(\gamma_1x)$, whenever the
notation denotes the corresponding formal tail product, has J-fraction
\[
 b_k=\gamma_{2k+1}+\gamma_{2k+2},\qquad
 \lambda_k=\gamma_{2k}\gamma_{2k+1}.
\]
For a numerical nonzero $\gamma_1$ this is ordinary division.
\item If $F=\Jfrac(b,\lambda;x)$, then
$(1-cx)^{-1}F(x/(1-cx))=\Jfrac(b+c,\lambda;x)$.
\end{enumerate}
\end{lemma}
\begin{proof}
Let $S_j=(1-\gamma_{j+1}xS_{j+1})^{-1}$ be the S-fraction tails.
For $V_j=S_jS_{j+1}$, direct multiplication gives
\[
 V_j^{-1}
 =1-(\gamma_{j+1}+\gamma_{j+2})x
   -\gamma_{j+2}\gamma_{j+3}x^2V_{j+2}.
\]
Indeed $S_{j+2}-1=\gamma_{j+3}xS_{j+2}S_{j+3}$.
Also $(S_0-1)/(\gamma_1x)=S_0S_1=V_0$, proving (ii).
For (i), rewrite
$S_0^{-1}=1-\gamma_1x-\gamma_1\gamma_2x^2V_1$ and use the same recursion.
For (iii), substitute $x/(1-cx)$ in a finite J-fraction and successively
multiply each denominator by $1-cx$.  The coefficient of every fixed power
of $x$ then stabilizes as the finite depth increases.
\end{proof}

\subsection{A self-contained route to the tangent fraction}

\begin{lemma}\label{pou:lem:tangentS}
The tangent moment generating series satisfies
\[
 \sum_{n\ge0}T_nx^n=\Sfrac(\gamma;x),\qquad \gamma_j=j(j+1).
\]
\end{lemma}
\begin{proof}
Put $f_k(z)=\tan^kz\,\sec^2z$, for $k\ge0$.
A differentiation and the identity $\sec^2z=1+\tan^2z$ give (with the
first term omitted when $k=0$)
\[
 Df_k=kf_{k-1}+(k+2)f_{k+1},\qquad f_k(0)=\ind_{k=0}.
\]
Thus $D^{2n}f_0(0)=T_n$ is a sum over walks of length $2n$ that start and
end at height zero, stay nonnegative, and have steps $\pm1$.
An upward crossing from height $j-1$ has weight $j+1$, and its downward
crossing has weight $j$.  A returning walk crosses each level equally often
in the two directions.  Its weight therefore equals that of a Dyck path
whose up steps have weight one and whose down step from height $j$ has
weight $j(j+1)$.  First-return decomposition at each height gives
$S_j=1+\gamma_{j+1}xS_{j+1}S_j$, the defining recursion of the S-fraction.
\end{proof}

\begin{theorem}[Bala's displayed conjecture]\label{pou:thm:bala}
Let $A(x)=\sum_{n\ge0}a_nx^n$.  Define
\begin{equation}\label{pou:eq:beta}
 \beta_{2j-1}=(2j+1)^2,\qquad
 \beta_{2j}=4j(j+1)\qquad(j\ge1).
\end{equation}
Then, in $\Z[[x]]$,
\begin{equation}\label{pou:eq:balafinal}
 A(x)=\frac{1+2x\Sfrac(\beta;x)}{1-x}.
\end{equation}
The successive weights inside $\Sfrac(\beta;x)$ are
$9,8,25,24,49,48,81,80,\ldots$.
\end{theorem}
\begin{proof}
Let $T(x)=\sum T_nx^n$ and
$V(x)=(T(x)-1)/(2x)=\sum T_{n+1}x^n/2$.
Lemma~\ref{pou:lem:contractions}(ii) applied to Lemma~\ref{pou:lem:tangentS}
gives the J-fraction of $V$:
\begin{equation}\label{pou:eq:VJ}
 b_k=8(k+1)^2,\qquad
 \lambda_k=4k(k+1)(2k+1)^2.
\end{equation}
The binomial transform identity is
$A(x)=(1-x)^{-1}T(x/(1-x))$.  Consequently
\[
 B(x):=\frac{(1-x)A(x)-1}{2x}
      =\frac1{1-x}V\left(\frac{x}{1-x}\right).
\]
Its J-diagonals are $8(k+1)^2+1$, with the same $\lambda_k$.
On the other hand, the even contraction of \eqref{pou:eq:beta} has
$b_0=9$ and, for $k\ge1$,
\[
 \beta_{2k}+\beta_{2k+1}=8(k+1)^2+1,\qquad
 \beta_{2k-1}\beta_{2k}=4k(k+1)(2k+1)^2.
\]
Thus the two J-fractions coincide coefficientwise.  Hence
$B=\Sfrac(\beta;x)$, which rearranges to \eqref{pou:eq:balafinal}.
\end{proof}

\begin{remark}
The result is an application of classical contraction and binomial-shift
identities, not a new theory of continued fractions.  Its value here is that
it proves exactly the later conjectural formula, with no finite-depth
assumption or experimental extrapolation.
\end{remark}

\section{The valuation observation and finite modular structure}\label{pou:sec:modular}

Even contraction of the tangent fraction and translation by $c$ give
\begin{equation}\label{pou:eq:AJ}
 \sum_{n\ge0}A_n(c)x^n=\Jfrac(b,\lambda;x),\qquad
 b_k=8k^2+8k+2+c,\quad \lambda_k=4k^2(4k^2-1).
\end{equation}
This also identifies $A_n(c)$ as weighted Motzkin path sums: horizontal
steps at height $k$ have weight $b_k$, upward steps weight one, and downward
steps from height $k$ weight $\lambda_k$.  First-return decomposition proves
this interpretation directly.

\begin{theorem}[An exact congruence]\label{pou:thm:mod9}
For every $n\ge1$, $a_n\equiv3\pmod9$.  In particular,
$v_3(a_n)=1$ for every $n\ge1$.
\end{theorem}
\begin{proof}
At $c=1$ the relevant initial weights are
\[
 b_0=3,\quad b_1=19,\quad \lambda_1=12,\quad \lambda_2=240.
\]
Every Motzkin path reaching height two contains the factor
$\lambda_1\lambda_2=2880$, divisible by nine.
Every path with two or more excursions from height zero contains
$\lambda_1^2$, also divisible by nine.  A path with one such excursion and
at least one horizontal step at height zero has a factor
$\lambda_1b_0$, again divisible by nine.
For $n\ge2$ the path consisting only of zero-height horizontal steps has
weight $3^n\equiv0\pmod9$.  The only remaining type is a single upward
step, then $n-2$ horizontal steps at height one, then a downward step.
Its weight is $12\cdot19^{n-2}\equiv3\pmod9$.
For $n=1$ the weight is $b_0=3$.  This proves the congruence at every index.
\end{proof}

\begin{corollary}\label{pou:cor:modrational}
For every positive integer $m$ and integer $c$, the formal series
$\sum A_n(c)x^n$ is rational over $\Z/m\Z$.  Its coefficient sequence
modulo $m$ is eventually periodic.
\end{corollary}
\begin{proof}
The weight $\lambda_m=4m^2(4m^2-1)$ vanishes modulo $m$.
Paths reaching height $m$ therefore vanish, so the series is the finite
J-fraction with heights $0,\ldots,m-1$.  Its denominator has constant term
one.  The coefficients obey a finite linear recurrence over a finite ring;
the finite vector of successive coefficients evolves deterministically
through a finite state space and is eventually periodic.
\end{proof}

The last corollary does not determine a minimal period.  It is a structural
consequence that can support later, more refined congruence calculations.
For the particular modulus nine, Theorem~\ref{pou:thm:mod9} says simply
$A(x)\equiv1+3x/(1-x)\pmod9$.

\section{Positive measures, uniqueness, and Hankel determinants}\label{pou:sec:moments}

\begin{proposition}[Exact moment representation]\label{pou:prop:moment}
For every integer $n\ge0$ and real $c\ge0$,
\begin{align}
 T_n&=\int_0^\infty y^{2n}\frac{y}{\sinh(\alpha y)}\dd y,
 \label{pou:eq:Tmoment}\\
 A_n(c)&=\int_0^\infty(c+y^2)^n\frac{y}{\sinh(\alpha y)}\dd y.
 \label{pou:eq:Amoment}
\end{align}
In particular, $y/\sinh(\alpha y)$ is a probability density on $[0,\infty)$.
The corresponding Stieltjes measure for $A_n(c)$ is
\begin{equation}\label{pou:eq:pushforward}
 \dd\mu_c(t)=\frac{\ind_{t>c}}{2\sinh(\alpha\sqrt{t-c})}\dd t.
\end{equation}
\end{proposition}
\begin{proof}
The classical cosine product \cite{dlmfcos} and its logarithmic derivative
give, for $|z|<\alpha$,
\[
 \tan z=\sum_{\substack{q\ge1\\q\ \mathrm{odd}}}
 \frac{2z}{\alpha^2q^2-z^2}.
\]
Comparison of Taylor coefficients yields
\begin{equation}\label{pou:eq:Tdirichlet}
 T_n=2\Gamma(2n+2)\alpha^{-2n-2}
       \sum_{q\ \mathrm{odd}\ge1}q^{-2n-2}.
\end{equation}
On the positive real axis,
$1/\sinh(\alpha y)=2\sum_{q\ \mathrm{odd}\ge1}e^{-\alpha qy}$.
Termwise integration is justified by nonnegativity and gives
\eqref{pou:eq:Tmoment}.  The binomial theorem gives \eqref{pou:eq:Amoment}.
The case $n=0$ has value $T_0=1$.  Finally, set $t=c+y^2$.
\end{proof}

The endpoint density in \eqref{pou:eq:pushforward} has an integrable
$(t-c)^{-1/2}$ singularity.  Its large-$t$ tail is exponential in $\sqrt t$,
not in $t$; this is why the ordinary moment generating series has zero
radius but a useful positive integral representation.

\begin{proposition}[Moment determinacy]\label{pou:prop:determinate}
For fixed $c\ge0$, the sequence $(A_n(c))_{n\ge0}$ has exactly one
representing probability measure on $[0,\infty)$.
\end{proposition}
\begin{proof}
There are constants $C,D>0$, depending on $c$, such that
$A_n(c)\le C D^{2n}(2n+1)!$.  For example,
$(c+y^2)^n\le2^n(c^n+y^{2n})$ and \eqref{pou:eq:Tdirichlet} give such a bound.
For any representing measure $\nu$, replace $t$ by the symmetric pair
$\pm\sqrt t$ to obtain a symmetric probability measure $\widetilde\nu$
on $\R$.  Its even moments are $A_n(c)$ and its odd moments are zero.
For sufficiently small $\delta>0$, monotone convergence gives
\[
 \int\cosh(\delta x)\dd\widetilde\nu(x)
 =\sum_{n\ge0}\frac{\delta^{2n}A_n(c)}{(2n)!}<\infty.
\]
Therefore its moment generating function is analytic in a neighborhood of
zero and is determined by its moments.  Two such symmetric measures have
the same moment generating function there, hence the same characteristic
function and the same measure.  Pushing forward by $x\mapsto x^2$ proves
uniqueness of $\nu$.
\end{proof}

\begin{proposition}[Hankel products]\label{pou:prop:hankel}
For $N\ge1$,
\begin{equation}\label{pou:eq:hankelA}
 \det\bigl(A_{i+j}(c)\bigr)_{0\le i,j<N}
 =\prod_{k=1}^{N-1}\bigl[4k^2(4k^2-1)\bigr]^{N-k}.
\end{equation}
If
$B_n(c)=\frac12\sum_{k=0}^n\binom nk c^{n-k}T_{k+1}$, then
\begin{equation}\label{pou:eq:hankelB}
 \det\bigl(B_{i+j}(c)\bigr)_{0\le i,j<N}
 =\prod_{k=1}^{N-1}\bigl[4k(k+1)(2k+1)^2\bigr]^{N-k}.
\end{equation}
Both determinants are independent of $c$ and strictly positive.
\end{proposition}
\begin{proof}
Here is an algebraic justification of the usual J-fraction determinant
rule.  On the monic polynomial basis $p_0=1,p_1,p_2,\ldots$, let
multiplication by the variable have recurrence
\[
 t p_k=p_{k+1}+b_kp_k+\lambda_kp_{k-1},\qquad \lambda_0p_{-1}=0.
\]
Declare this basis orthogonal, with squared norms
$h_0=1$ and $h_k=\lambda_1\cdots\lambda_k$.
The recurrence makes multiplication by $t$ symmetric for this bilinear
form.  Its moments $\langle1,t^n1\rangle$ are exactly the weighted Motzkin
path sums and hence the coefficients of the J-fraction.
The change from monomials to monic $p_k$ is unit triangular, so the Hankel
determinant is $\prod_{k=0}^{N-1}h_k$.
Apply this to \eqref{pou:eq:AJ} and to the translated version of
\eqref{pou:eq:VJ}.  The formulas follow without any analytic convergence
assumption on the generating functions.
\end{proof}

\section{Exact exponential sectors and all-order fixed-shift expansions}
\label{pou:sec:sectors}

Let
\begin{equation}\label{pou:eq:FandH}
 F_n=2\Gamma(2n+2)\alpha^{-2n-2},\qquad
 H_n(z)=\sum_{j=0}^n\frac{(z/4)^j}{j!\fall{n+\tfrac12}{j}},
\end{equation}
where $\fall u0=1$.  The smallest factor in a nonempty denominator of
$H_n$ is at least $3/2$, so no zero denominator occurs.

\begin{theorem}[Exact sectors and an explicit tail bound]\label{pou:thm:sectors}
For every $n\ge0$ and $c\ge0$,
\begin{equation}\label{pou:eq:exactsectors}
 \frac{A_n(c)}{F_n}
 =\sum_{q\ \mathrm{odd}\ge1}q^{-2n-2}H_n(c\alpha^2q^2).
\end{equation}
The series is convergent and all its terms are nonnegative.
If $Q\ge1$ is odd and $R_{n,Q}(c)$ is the contribution to $A_n(c)$
from $q\ge Q$, then
\begin{equation}\label{pou:eq:tailbound}
 0\le\frac{R_{n,Q}(c)}{F_n}
 \le \left(1+\frac Q2\right)
 Q^{-2n-2}H_n(c\alpha^2Q^2).
\end{equation}
\end{theorem}
\begin{proof}
Insert \eqref{pou:eq:Tdirichlet} into the binomial transform and put $j=n-k$.
The duplication identity, or cancellation of consecutive factorials, gives
\[
 \binom nj\frac{\Gamma(2n-2j+2)}{\Gamma(2n+2)}
 =\frac1{4^j j!\fall{n+\tfrac12}{j}}.
\]
After factoring $q^{-2n-2}$, this is exactly \eqref{pou:eq:exactsectors}.
At fixed $n$ the last polynomial term decays as $q^{-2}$; every other term
decays faster, proving convergence.
For every real $r\ge2$, monotonicity and an integral comparison give
\[
 \sum_{j\ge0}(Q+2j)^{-r}
 \le Q^{-r}+\frac12\int_Q^\infty u^{-r}\dd u
 =Q^{-r}\left(1+\frac{Q}{2(r-1)}\right)
 \le Q^{-r}(1+Q/2).
\]
Apply this separately to each nonnegative coefficient of $H_n$; its
corresponding exponent is $r=2n+2-2j\ge2$.
\end{proof}

\begin{theorem}[A coefficient algorithm to every algebraic order]
\label{pou:thm:Hallorders}
Uniformly for $z$ in each fixed compact subset of $\mathbb C$,
\begin{equation}\label{pou:eq:Hexpansion}
 H_n(z)\sim\sum_{r\ge0}\frac{C_r(z)}{n^r},
\end{equation}
where
\begin{equation}\label{pou:eq:Cformula}
 C_r(z)=\sum_{j=0}^r\frac{(z/4)^j}{j!}
 [t^{r-j}]\prod_{\ell=0}^{j-1}
       \bigl(1+(\tfrac12-\ell)t\bigr)^{-1}.
\end{equation}
In particular,
\begin{align*}
 C_0(z)&=1,& C_1(z)&=\frac z4,\\
 C_2(z)&=-\frac z8+\frac{z^2}{32},&
 C_3(z)&=\frac z{16}+\frac{z^3}{384},\\
 C_4(z)&=-\frac z{32}+\frac{z^2}{128}
          +\frac{z^3}{256}+\frac{z^4}{6144}.
\end{align*}
\end{theorem}
\begin{proof}
For a fixed $j$ the summand is
\[
 \frac{(z/4)^j}{j!n^j}
 \prod_{\ell=0}^{j-1}
       \left(1+\frac{1/2-\ell}{n}\right)^{-1}.
\]
Taylor expansion of the finite product gives \eqref{pou:eq:Cformula}.
It remains to justify discarding unbounded $j$.
For $j\le n/2$, every falling-factorial factor is at least $n/2$, so the
absolute value of the summand is at most $(|z|/(2n))^j/j!$.
The sum of these bounds over $j\ge M$ is $O(n^{-M})$, uniformly for bounded
$z$.  For $j>n/2$, at least $\lfloor n/2\rfloor$ factors are at least
$n/2$ and all remaining factors exceed one.  Summing the numerator divided
by $j!$ bounds this part by
$e^{|z|/4}(n/2)^{-\lfloor n/2\rfloor}$.
The finite terms with $j<M$ have uniform Taylor remainders, completing the
proof at each prescribed order.
\end{proof}

Combining these two theorems gives an explicit description of every fixed
odd-pole sector.  For a fixed odd cutoff $Q$ and fixed $c$, the omitted
sectors are $O(F_nQ^{-2n-2})$.  This statement uses the exact $H_n$ at the
cutoff, not a claim uniform in arbitrarily growing $Q$ after replacing
$H_n$ by a fixed-$z$ expansion.

\subsection{Why the exact sectors matter}

The first exponentially small correction is meaningful in the form
\begin{equation}\label{pou:eq:threecorrection}
 \frac{A_n(c)}{F_nH_n(c\alpha^2)}
 =1+3^{-2n-2}\frac{H_n(9c\alpha^2)}{H_n(c\alpha^2)}
     +O_c(5^{-2n-2}).
\end{equation}
The exact denominator is essential.  Replacing it by a fixed finite
inverse-power approximation introduces a polynomial error that is much
larger than $3^{-2n}$.  Formula~\eqref{pou:eq:threecorrection} avoids that
error-order confusion and supplies a rigorous positive residual bound
through \eqref{pou:eq:tailbound}.

\subsection{The logarithmic expansion and the prior leading equivalent}

Write
$\log(\sum_{r\ge0}C_r(z)t^r)=\sum_{r\ge1}\ell_r(z)t^r$.
The first six logarithmic coefficients, generated by the supplied code,
are
\[
\begin{aligned}
\ell_{1}(z)&=\frac{z}{4}\\
\ell_{2}(z)&=- \frac{z}{8}\\
\ell_{3}(z)&=\frac{z^{2}}{32} + \frac{z}{16}\\
\ell_{4}(z)&=- \frac{z^{2}}{64} - \frac{z}{32}\\
\ell_{5}(z)&=\frac{z^{3}}{96} + \frac{z^{2}}{64} + \frac{z}{64}\\
\ell_{6}(z)&=\frac{z^{3}}{192} - \frac{z^{2}}{128} - \frac{z}{128}\\

\end{aligned}
\]
For fixed $c\ge0$, Stirling's expansion \cite{dlmfgamma} and the exponentially
small sectors give
\begin{equation}\label{pou:eq:logA}
 \log A_n(c)\sim
 2n\log\frac{2n}{\alpha e}+\frac32\log n
 +\log\frac{8\sqrt\pi}{\alpha^2}
 +\sum_{r\ge1}\frac{h_r(c)}{n^r}.
\end{equation}
The coefficient rule is
\begin{equation}\label{pou:eq:hgeneral}
 h_r(c)=\ell_r(c\alpha^2)+\frac{(-1)^{r+1}}{r2^r}
 +\begin{cases}
 \displaystyle\frac{B_{r+1}}{(r+1)r\,2^r},&r\text{ odd},\\[3pt]
 0,&r\text{ even},
 \end{cases}
\end{equation}
where $B_j$ are Bernoulli numbers.  For example,
\begin{align}
 h_1(c)&=\frac{13}{24}+\frac{c\alpha^2}{4},\nonumber\\
 h_2(c)&=-\frac18-\frac{c\alpha^2}{8},\label{pou:eq:firsth}\\
 h_3(c)&=\frac{119}{2880}+\frac{c\alpha^2}{16}
                       +\frac{c^2\alpha^4}{32}.\nonumber
\end{align}
One way to verify the constants is to use
$\Gamma(2n+2)=(2n+1)\Gamma(2n+1)$: the extra logarithm is
$\log(1+1/(2n))$, and the remaining correction is the ordinary Stirling
series at $2n$.

At $c=1$, exponentiating the leading part gives
\begin{equation}\label{pou:eq:knownleading}
 a_n\sim\frac{2^{4n+5}n^{2n+3/2}}{e^{2n}\pi^{2n+3/2}},
\end{equation}
which is the leading equivalent already recorded by Kotesovec in 2015
\cite{oeis125054}.  Formula~\eqref{pou:eq:knownleading} is therefore a
consistency check, not a newly solved asymptotic problem.  The all-order
and exponentially resolved refinements above go beyond that leading
formula in the inspected entry.

\section{Quadratic shifts: the geometry of a first-order transition}
\label{pou:sec:geometry}

The fixed-$c$ expansion cannot be inserted blindly when $c$ grows like
$n^2$.  In that regime an entirely different contribution competes with
the factorial tail.  Set $c=\tau n^2$ in \eqref{pou:eq:Amoment}, then put $y=nu$:
\begin{equation}\label{pou:eq:scaledintegral}
 A_n(\tau n^2)=n^{2n+2}\int_0^\infty
  (\tau+u^2)^n\frac{u}{\sinh(\alpha nu)}\dd u.
\end{equation}
Near a fixed positive $u$, the hyperbolic factor is asymptotic to
$2e^{-\alpha nu}$; the phase is
\begin{equation}\label{pou:eq:phase}
 \phi_\tau(u)=\log(\tau+u^2)-\alpha u.
\end{equation}
The region $y=O(1)$, equivalently $u=O(1/n)$, is an endpoint contribution
with leading value $(\tau n^2)^n$.  It must not be estimated by replacing
$1/\sinh(\alpha y)$ by its large-$y$ first exponential.

For $0<\tau<\alpha^{-2}$ the stationary points are
\[
 u_\pm=\frac{1\pm\sqrt{1-\alpha^2\tau}}\alpha.
\]
The larger one is a local maximum.  Parameterize it by
\begin{equation}\label{pou:eq:sparam}
 s=\alpha u_+=1+\sqrt{1-\alpha^2\tau},\qquad
 \tau=\frac{s(2-s)}{\alpha^2},\qquad 1<s<2.
\end{equation}
The saddle's logarithmic advantage over the endpoint is
\begin{equation}\label{pou:eq:Delta}
 \Delta(\tau)=\phi_\tau(u_+)-\log\tau
 =\log\frac2{2-s}-s.
\end{equation}
As a function of $s$, its derivative is $(s-1)/(2-s)>0$.
It is negative at $s=1$ and tends to infinity at $s=2$.
Thus there is a unique $s_*$ in $(1,2)$ solving \eqref{pou:eq:starintro}.

The curvature, amplitude, and derivative of the phase difference are
\begin{align}
 \kappa(\tau)&=-\phi_\tau''(u_+)
        =\alpha^2\frac{s-1}{s},\label{pou:eq:kap}\\
 K(\tau)&=2u_+\sqrt{\frac{2\pi}{\kappa(\tau)}}
        =\frac{2\sqrt{2\pi}s^{3/2}}{\alpha^2\sqrt{s-1}},\label{pou:eq:K}\\
 -\Delta'(\tau)&=\frac{\alpha^2}{2(2-s)}.\label{pou:eq:deltaprime}
\end{align}
In particular, with $d_*=-\Delta'(\tau_*)$ and $K_*=K(\tau_*)$,
\begin{align*}
 s_*&=1.5936242600400400923\ldots,\\
 \tau_*&=0.2624665433684190267\ldots,\\
 d_*&=3.0358617132452985767\ldots,\\
 K_*&=5.3052182371146210849\ldots.
\end{align*}
These displayed digits are numerical evaluations of exact definitions,
not interval-certified decimal bounds.

\writenote{Notation across reports: $s_*=2+W_0(-2e^{-2})=1.59362426\ldots$
is the constant written $\tau$ in the A277364 report \cite{bellreport};
here $\tau$ is the shift ratio $c/n^2$, and
$\tau_*=s_*(2-s_*)/\alpha^2=0.26246654\ldots$ is a different number. Do
not read the A277364 report's $\tau$ as this report's $\tau_*$. The
shared equation $2-s=2e^{-s}$ arises there from the omitted tail of a
Stirling sum and here from the endpoint--saddle balance $\Delta(\tau)=0$; neither
report uses the other's argument.}

\begin{proposition}[Fixed-parameter regimes]\label{pou:prop:fixedregimes}
If $0<\tau<\tau_*$ is fixed, then
\[
 A_n(\tau n^2)\sim n^{2n}e^{n\phi_\tau(u_+)}K(\tau)n^{3/2}.
\]
If $\tau>\tau_*$ is fixed, then
$A_n(\tau n^2)\sim(\tau n^2)^n$.
At $\tau=\tau_*$ the saddle, not the endpoint, still dominates by the
factor $K_*n^{3/2}$.
\end{proposition}
\begin{proof}
When the larger stationary point exists, comparison of the endpoint value
$\log\tau$ and the saddle value is exactly the sign of $\Delta$.
There are no other local maxima.  For $\tau\ge\alpha^{-2}$ the phase has
no larger interior value, so the endpoint is the global maximum.
The endpoint integral tends to one after division by $(\tau n^2)^n$;
the saddle has amplitude $2n^2u_+$ and Gaussian width $n^{-1/2}$ in $u$,
giving the factor in \eqref{pou:eq:K}.  The localization argument in the next
section supplies uniform errors and proves these leading assertions.
\end{proof}

The terminology ``first-order'' refers to the discontinuous limiting
fraction of the binomial sum's active index, proved in
Section~\ref{pou:sec:probability}.  It is not merely a change in a prefactor.

\section{A uniform two-contribution expansion to all orders}\label{pou:sec:uniform}

Choose a sufficiently small closed interval $I$ about $\tau_*$ contained
in $(0,\alpha^{-2})$.  All constants implicit in the following theorem may
depend on $I$ and the fixed truncation order, but not on $\tau\in I$.
For a list $v_1,\ldots,v_r$, let $e_j(v_1,\ldots,v_r)$ denote its elementary
symmetric polynomial, with $e_0=1$ and $e_j=0$ outside its natural range.

\begin{theorem}[Uniform all-order endpoint--saddle expansion]
\label{pou:thm:uniform}
For every fixed integer $M\ge1$, uniformly for $\tau\in I$,
\begin{align}
 \frac{A_n(\tau n^2)}{(\tau n^2)^n}
 ={}&\sum_{m=0}^{M-1}\frac{E_m(\tau)}{n^m}
 +K(\tau)n^{3/2}e^{n\Delta(\tau)}
     \sum_{m=0}^{M-1}\frac{L_m(\tau)}{n^m}\nonumber\\
 &+O_{I,M}\!\left(n^{-M}
       \left[1+n^{3/2}e^{n\Delta(\tau)}\right]\right).
 \label{pou:eq:uniform}
\end{align}
Here $E_0=L_0=1$ and
\begin{equation}\label{pou:eq:Ecoeff}
 E_m(\tau)=\sum_{k=0}^{m}
 \frac{T_k}{k!\tau^k}(-1)^{m-k}
 e_{m-k}(0,1,\ldots,k-1).
\end{equation}
The saddle coefficients have the finite Gaussian coefficient rule
\begin{equation}\label{pou:eq:Lcoeff}
 L_m(\tau)=[\eps^{2m}]\E\left[
 \left(1+\frac{\eps Z}{u_+}\right)
 \exp\left\{\sum_{r\ge3}\frac{\phi_\tau^{(r)}(u_+)}{r!}
                  Z^r\eps^{r-2}\right\}\right],
\end{equation}
where $Z\sim\Normal(0,\kappa(\tau)^{-1})$ and the expectation is applied
to each polynomial coefficient.  In particular,
\begin{align}
 E_1(\tau)&=\frac2\tau,&E_2(\tau)&=\frac8{\tau^2},&
 E_3(\tau)&=\frac{136}{3\tau^3}-\frac8{\tau^2},\label{pou:eq:Es}\\
 L_1(\tau)&=\frac{2s^3+18s^2-60s+45}{24(s-1)^3}.
 \label{pou:eq:Lone}
\end{align}
\end{theorem}

\begin{proof}
We give the localization and remainder details, because the endpoint
and the interior saddle live on different scales.

\emph{Endpoint.}  Fix a sufficiently small $\eta>0$, uniformly over
$\tau\in I$, so that $\eta/\min I<\alpha/2$ and $\eta$ lies below the
smaller stationary point throughout $I$.
On $0\le y\le\eta n$,
\[
 \left(1+\frac{y^2}{\tau n^2}\right)^n
 \le\exp\left(\frac{y^2}{\tau n}\right)
 \le\exp\left(\frac{\eta y}{\min I}\right).
\]
Multiplication by $y/\sinh(\alpha y)$ therefore leaves an integrable
exponential majorant, uniformly in $n$ and $\tau$.
At each prescribed order, expand
\[
 \left(1+\frac{y^2}{\tau n^2}\right)^n
 =\sum_{k=0}^n\frac{\fall nk}{k!}
                \frac{y^{2k}}{\tau^k n^{2k}}.
\]
For a fixed coefficient of $n^{-m}$, only $k\le m$ can contribute;
expanding $\fall nk=n^k\prod_{j=0}^{k-1}(1-j/n)$ gives precisely
\eqref{pou:eq:Ecoeff} after integration.

For a fully uniform remainder bound, introduce the analytic function
\[
 G_y(t)=\exp\!\left(\frac{\log(1+y^2t^2/\tau)}t\right),\qquad G_y(0)=1.
\]
Choose $\rho>0$ so small that $\rho^2/\min I<1/2$ and
$2\rho/\min I<\alpha/2$.  On the complex circle
$|t|=\rho/(1+y)$ the logarithm is its branch near one and
$|G_y(t)|\le\exp(2\rho y/\min I)$.
Reduce $\eta$ to at most $\rho/4$.  For all sufficiently large $n$ and
$0\le y\le\eta n$, the evaluation point $t=1/n$ is inside half this
circle.  Cauchy's coefficient estimate bounds the Taylor remainder after
$t^{M-1}$ by
\[
 2\exp(2\rho y/\min I)
     \left(\frac{1+y}{\rho n}\right)^M.
\]
Its product with $y/\sinh(\alpha y)$ is integrable uniformly in $n$ and
$\tau$, proving the required $O(n^{-M})$ integrated error.
The Taylor coefficients agree with the finite formal calculation above.
Extending their integrals from $\eta n$ to infinity changes only an
exponentially small amount.

\emph{Interior saddle.}  Choose a small neighborhood of $u_+(\tau)$
whose distance from zero is uniformly positive and whose curvature is
bounded away from zero.  In this region
\[
 \frac1{\sinh(\alpha nu)}
 =2e^{-\alpha nu}\bigl(1+O(e^{-\rho n})\bigr)
\]
for some $\rho>0$.  Set $u=u_++n^{-1/2}z$ in the first exponential term
of \eqref{pou:eq:scaledintegral}.  The quadratic term gives the centered
Gaussian with variance $\kappa^{-1}$, and the amplitude $u/u_+$ is
$1+n^{-1/2}z/u_+$.  Taylor's formula for the remaining phase gives
\eqref{pou:eq:Lcoeff}.  Every coefficient involves only finitely many
derivatives and Gaussian moments.  Odd powers of $n^{-1/2}$ integrate to
zero.  The standard localized Laplace estimate, obtained by restricting
first to $|z|\le n^\delta$ and then using Gaussian domination on its
complement, gives relative error $O(n^{-M})$ uniformly on $I$
\cite{dlmflaplace}.  The prefactor is $K(\tau)n^{3/2}e^{n\Delta(\tau)}$.

\emph{Complement.}  Outside the endpoint region and the saddle
neighborhood, $\phi_\tau(u)$ is less than
$\max\{\log\tau,\phi_\tau(u_+)\}$ by a fixed positive amount on each
sufficiently large compact $u$ interval.  For large $u$, the term
$-\alpha u$ dominates $\log(\tau+u^2)$ uniformly, giving the same type of
exponential bound on the unbounded tail.  The compact interval $I$ may be
reduced so that these gaps remain uniform.
Thus this complement is exponentially small relative to the larger of
the two contributions, and is absorbed by the remainder in
\eqref{pou:eq:uniform}.

Finally, evaluating \eqref{pou:eq:Ecoeff} yields \eqref{pou:eq:Es}.
For the first saddle correction, Gaussian moments give
\[
 L_1=\frac{\phi'''(u_+)}{2u_+\kappa^2}
       +\frac{\phi''''(u_+)}{8\kappa^2}
       +\frac{5\phi'''(u_+)^2}{24\kappa^3}.
\]
Substitution of \eqref{pou:eq:sparam} and \eqref{pou:eq:kap} simplifies to
\eqref{pou:eq:Lone}.
\end{proof}

\begin{remark}[Scope of the uniform expansion]
This theorem is uniform near the coexistence point, not near the
coalescing stationary points at $\tau=\alpha^{-2}$.
It also does not make a fixed-$c$ expansion uniform up to $c\asymp n^2$.
The two asymptotic problems use different localizations.  If one of the
terms in \eqref{pou:eq:uniform} is exponentially smaller than the other, a
finite inverse-power error in the larger term can mask it; the underlying
localized integrals still have their separate expansions.
\end{remark}

\section{The moving window and the probability of a macroscopic index}
\label{pou:sec:probability}

\begin{theorem}[Finite-size coexistence window]\label{pou:thm:window}
For fixed bounded real $w$, define
\begin{equation}\label{pou:eq:window}
 \tau_n(w)=\tau_*+
 \frac{\tfrac32\log n+\log K_*+w}{d_*n}.
\end{equation}
Then, uniformly for $w$ in compact real sets,
\begin{equation}\label{pou:eq:windowlimit}
 \frac{A_n(\tau_n(w)n^2)}{(\tau_n(w)n^2)^n}
 =1+e^{-w}+O\!\left(\frac{(\log n)^2}{n}\right).
\end{equation}
\end{theorem}
\begin{proof}
Write $\delta_n=\tau_n-\tau_*$.
Taylor expansion gives
$n\Delta(\tau_n)=-d_*n\delta_n+O(n\delta_n^2)$ and
$\log K(\tau_n)=\log K_*+O(\delta_n)$.
Since $\delta_n=O((\log n)/n)$,
\[
 \log\bigl(K(\tau_n)n^{3/2}e^{n\Delta(\tau_n)}\bigr)
 =-w+O((\log n)^2/n).
\]
The endpoint factor is $1+O(1/n)$ and the saddle relative correction is
$1+O(1/n)$ by Theorem~\ref{pou:thm:uniform}.  Combining them proves the result.
\end{proof}

A useful more accurate centering is the unique nearby solution of
\begin{equation}\label{pou:eq:exactbalance}
 n\Delta(\tau)+\tfrac32\log n+\log K(\tau)=-w.
\end{equation}
Its derivative with respect to $\tau$ is $-d_*n+O(\log n)$, so existence
and uniqueness in the relevant neighborhood follow for large $n$.
Under this centering the leading saddle factor is exactly $e^{-w}$ and
the error in \eqref{pou:eq:windowlimit} improves to $O(1/n)$.
Using the coefficients $E_m,L_m$ in the balance equation gives successive
higher-order corrections.  Equation~\eqref{pou:eq:window} is the simple explicit
first approximation, not the most accurate numerical prescription.

\subsection{A positive random-index representation}

For $c>0$ define a random index $K_n\in\{0,\ldots,n\}$ by
\begin{equation}\label{pou:eq:Klaw}
 \Prob_c(K_n=k)=\frac{\binom nk c^{n-k}T_k}{A_n(c)}.
\end{equation}
Under the probability measure on $y\ge0$ proportional to
$(c+y^2)^n y/\sinh(\alpha y)$, conditionally take
\begin{equation}\label{pou:eq:conditionalbinomial}
 K_n\mid y\ \sim\ \operatorname{Bin}
 \left(n,\frac{y^2}{c+y^2}\right).
\end{equation}
Integrating the conditional mass function recovers \eqref{pou:eq:Klaw}, so
this is an exact mixture construction, not an approximation.

\begin{theorem}[Discontinuous fraction and critical mixture]
\label{pou:thm:mixture}
If $0<\tau<\tau_*$ is fixed and $c=\tau n^2$, then
$K_n/n\to s(\tau)/2$ in probability.
If $\tau>\tau_*$ is fixed, then $K_n\to0$ in probability.
In the moving window \eqref{pou:eq:window},
\begin{equation}\label{pou:eq:twopoint}
 \frac{K_n}{n}\ \Longrightarrow\
 \frac1{1+e^{-w}}\delta_0+
 \frac{e^{-w}}{1+e^{-w}}\delta_{s_*/2}.
\end{equation}
The macroscopic location at coexistence is
$s_*/2=0.796812130020020\ldots$.
\end{theorem}
\begin{proof}
The endpoint tilted measure is concentrated on $y=O(1)$, while the saddle
measure has $y/n\to u_+$.  Conditional binomial variance divided by $n^2$
is at most $1/(4n)$.
The conditional mean fraction tends to zero in the endpoint region and,
in the saddle region, to
\[
 \frac{u_+^2}{\tau+u_+^2}=\frac s2.
\]
The two relative masses follow from \eqref{pou:eq:uniform} and
\eqref{pou:eq:windowlimit}.  In the fixed endpoint-dominated regime the stronger
claim follows directly from
$\Prob(K_n=0)=c^n/A_n(c)\to1$.
\end{proof}

\begin{theorem}[Gaussian fluctuations within the macroscopic component]
\label{pou:thm:clt}
For fixed $0<\tau<\tau_*$,
\begin{equation}\label{pou:eq:clt}
 \frac{K_n-ns(\tau)/2}{\sqrt n}
 \ \Longrightarrow\ \Normal\left(0,
   v(\tau)\right),\qquad
 v(\tau)=\frac{s(2-s)}{4(s-1)}.
\end{equation}
In the critical window the same conclusion holds conditionally on the
saddle component, with $s=s(\tau_n)$ and variance tending to $v(\tau_*)$.
Equivalently, one may condition on $K_n/n>\eps$ for any fixed
$0<\eps<s_*/2$.
\end{theorem}
\begin{proof}
Localized Laplace expansion gives
$\sqrt n(y/n-u_+)\Longrightarrow\Normal(0,\kappa^{-1})$.
Let $p(u)=u^2/(\tau+u^2)$.  Uniformly near $u_+$, the conditional
binomial characteristic function for its centered fluctuations converges
to that of $\Normal(0,p(u_+)(1-p(u_+)))$.
At the $\sqrt n$ scale the variation in its conditional mean is
$p'(u_+)\sqrt n(y/n-u_+)$.  The limiting characteristic function is thus
the product of the two Gaussian factors, with variance
\[
 p(u_+)(1-p(u_+))+\frac{p'(u_+)^2}{\kappa}.
\]
Since $p(u_+)=s/2$ and $p'(u_+)=\alpha(2-s)/2$, simplification gives
\eqref{pou:eq:clt}.  Endpoint leakage under the stated conditioning tends to
zero by the separated concentration regions.  The same localization and
conditional characteristic-function argument is uniform for
$\tau_n\to\tau_*$.
\end{proof}

The critical limiting law is a two-point law at scale $n$, with a Gaussian
width of order $\sqrt n$ around its nonzero point.  A single Gaussian
approximation to the entire distribution would therefore be incorrect.

\section{A universality theorem for exponentially decaying moment tails}
\label{pou:sec:universal}

In this section $\alpha$ is an arbitrary positive tail rate; in
Section~\ref{pou:sec:inverse} it again denotes $\pi/2$.
The logarithmic displacement is not specific to tangent numbers.  The
following result identifies the parameters on which it depends.
It is a consequence of a competition between an unscaled probability-mass
contribution and a scaled interior Laplace saddle; it does not assert a new
general method of asymptotic integration.

\begin{theorem}[Universal shifted-moment transition]\label{pou:thm:universal}
Let $\mu$ be a probability measure on $[0,\infty)$ whose restriction beyond
some finite point has a density
\begin{equation}\label{pou:eq:generaltail}
 w(y)=Cy^\beta e^{-\alpha y}\bigl(1+O(y^{-1})\bigr),
 \qquad C,\alpha>0,\quad \beta\in\R.
\end{equation}
Fix $p>1$ and put
\[
 m_k=\int y^{pk}\dd\mu(y),\qquad
 M_n(c)=\sum_{k=0}^n\binom nk c^{n-k}m_k
       =\int(c+y^p)^n\dd\mu(y).
\]
There is a unique $s_p\in(p-1,p)$ such that
\begin{equation}\label{pou:eq:generalstar}
 p-s_p=pe^{-s_p},\qquad
 \tau_{p,*}=\frac{s_p^{p-1}(p-s_p)}{\alpha^p}.
\end{equation}
Define
\begin{align}
 \kappa_*&=\alpha^2\frac{s_p-p+1}{s_p},&
 K_{p,*}&=C\left(\frac{s_p}{\alpha}\right)^\beta
                      \sqrt{\frac{2\pi}{\kappa_*}},\label{pou:eq:generalK}\\
 d_{p,*}&=\frac{\alpha^p}{p s_p^{p-2}(p-s_p)},&
 b&=\beta+\frac12.\label{pou:eq:generald}
\end{align}
For
\begin{equation}\label{pou:eq:generalwindow}
 \tau_n(w)=\tau_{p,*}+
       \frac{b\log n+\log K_{p,*}+w}{d_{p,*}n},
\end{equation}
we have, uniformly for bounded real $w$,
\begin{equation}\label{pou:eq:generallimit}
 \frac{M_n(\tau_n(w)n^p)}{(\tau_n(w)n^p)^n}
 =1+e^{-w}
 +O\left(n^{1-p}+\frac{(\log n)^2}{n}\right).
\end{equation}
The random binomial index with masses proportional to
$\binom nk c^{n-k}m_k$ satisfies
\begin{equation}\label{pou:eq:generalmixture}
 K_n/n\Longrightarrow
 \frac1{1+e^{-w}}\delta_0+
 \frac{e^{-w}}{1+e^{-w}}\delta_{s_p/p}.
\end{equation}
For fixed $0<\tau<\tau_{p,*}$, the centered index
$(K_n-ns/p)/\sqrt n$ converges to a centered Gaussian with variance
\begin{equation}\label{pou:eq:generalvariance}
 v_p(\tau)=\frac{s(p-s)}{p^2(s-p+1)},
 \qquad \tau=\frac{s^{p-1}(p-s)}{\alpha^p},\quad p-1<s<p.
\end{equation}
The same variance, evaluated at $s_p$, describes the macroscopic component
at coexistence.
\end{theorem}

\begin{proof}
The relevant phase after $y=nu$ is
\[
 \phi_{\tau,p}(u)=\log(\tau+u^p)-\alpha u.
\]
Its positive stationary points satisfy
\[
 \tau=\frac{s^{p-1}(p-s)}{\alpha^p},\qquad s=\alpha u.
\]
The larger stationary point has $p-1<s<p$ and negative second derivative
$-\alpha^2(s-p+1)/s$.  There are at most two positive stationary points,
because $s^{p-1}(p-s)$ increases up to $p-1$ and then decreases to zero.
At the larger point the excess over the endpoint phase is
\begin{equation}\label{pou:eq:generalDelta}
 \Delta_p(\tau)=\log\frac p{p-s}-s.
\end{equation}
As a function of $s\in(p-1,p)$, its derivative is
$(s-p+1)/(p-s)>0$.  Its value at the left endpoint is
$\log p-(p-1)<0$, while it tends to infinity at the right endpoint.
This proves existence and uniqueness in \eqref{pou:eq:generalstar}.
Differentiating the maximized phase with respect to $\tau$ gives
$-\Delta_p'(\tau)=\alpha^p/[p s^{p-2}(p-s)]$.

Choose a small endpoint region $y\le\eta n$, uniformly near the transition.
For $p>1$,
\[
 \left(1+\frac{y^p}{\tau n^p}\right)^n
 \le\exp\left(\frac{\eta^{p-1}y}{\tau}\right).
\]
Taking $\eta$ small enough preserves an integrable exponential majorant
against the tail of $\mu$.  Moreover $e^v-1\le ve^v$, so integration of
this bound shows that the normalized endpoint contribution is
$1+O(n^{1-p})$.  Compactly supported mass, including atoms, contributes
in the same way and is included in this leading one.

Near the larger saddle, the tail assumption gives the amplitude
$C n^{\beta+1}u^\beta$ times $1+O(1/n)$, uniformly in a compact positive
$u$ interval.  The Gaussian width is $n^{-1/2}$, so the normalized saddle
contribution is
\[
 K_p(\tau)n^{\beta+1/2}e^{n\Delta_p(\tau)}(1+O(1/n)),
 \quad
 K_p(\tau)=C(s/\alpha)^\beta
 \sqrt{\frac{2\pi s}{\alpha^2(s-p+1)}}.
\]
The $O(y^{-1})$ error need not be differentiable: a uniform relative bound
on the localized integral suffices.  The smooth main amplitude has the
usual Laplace estimate.  The phase gap on the remaining region is handled
as in Theorem~\ref{pou:thm:uniform}.  Inserting \eqref{pou:eq:generalwindow} and
Taylor expanding $\Delta_p$ and $\log K_p$ proves
\eqref{pou:eq:generallimit}.  The error $n^{1-p}$ must be retained when
$1<p<2$; replacing it by $1/n$ would in general be false.

The exact conditional construction is now
$K_n\mid y\sim\operatorname{Bin}(n,y^p/(c+y^p))$.
At the saddle the success probability is $s/p$, so the same concentration
argument proves \eqref{pou:eq:generalmixture}.
For the Gaussian fluctuations,
\[
 p_0=\frac sp,\qquad
 p'_0=\frac{\alpha(p-s)}p,\qquad
 \kappa=\alpha^2\frac{s-p+1}{s}.
\]
Their variance is $p_0(1-p_0)+(p'_0)^2/\kappa$, which equals
\eqref{pou:eq:generalvariance}.  The conditional characteristic-function proof
used in Theorem~\ref{pou:thm:clt} applies unchanged.
\end{proof}

\subsection{What the universality theorem does and does not cover}

The transition location depends on $p$ and the exponential rate $\alpha$,
not on the compact part of the measure or the tail power $\beta$.
The displacement of the finite-size window does depend on $\beta$:
its logarithmic coefficient is $\beta+1/2$.
This coefficient may be negative, zero, or positive; the theorem requires
none of those cases to be excluded.
The assumption that $\mu$ is a probability measure fixes the endpoint
amplitude to one.  Without that normalization the mass of the measure
must be inserted explicitly in the balance equation.

The theorem gives a leading uniform transition and a bulk central limit
law.  It does not claim an all-order expansion for an arbitrary
$O(y^{-1})$ tail perturbation.  An all-order tail expansion, together with
appropriate regularity or controlled remainders, would support an
all-order extension; the tangent case has that additional structure.
The boundary case $p=1$ is excluded because its interior stationary point
coalesces with the endpoint as the critical parameters are approached.

\subsection{Examples beyond A125054}

The tangent density has $C=2$, $\beta=1$ and $p=2$, recovering the
$\tfrac32\log n$ displacement.
The positive density $\sech(\pi y/2)$ on $[0,\infty)$ has even moments
the secant numbers $E_{2k}$, OEIS A000364 \cite{oeis364}.
For example, integrating its geometric exponential expansion gives
$2\Gamma(2k+1)\alpha^{-2k-1}\sum_{j\ge0}(-1)^j(2j+1)^{-2k-1}$,
the usual secant partial-fraction coefficient.  In this case $C=2$,
$\beta=0$, so the same quadratic-shift transition location has a
$\tfrac12\log n$ displacement instead.

For a gamma density
\[
 \dd\mu(y)=\frac{\alpha^b}{\Gamma(b)}y^{b-1}e^{-\alpha y}\dd y,
 \qquad b>0,
\]
the moments are explicitly
$m_k=\alpha^{-pk}\Gamma(pk+b)/\Gamma(b)$.
Here the logarithmic coefficient is $b-1/2$.
Thus the theorem covers factorial subsequences and generalized factorial
moment transforms as well as alternating-permutation numbers.

More generally, a tail $Dy^\gamma e^{-a y^r}$ can be reduced by the
substitution $v=y^r$ to the theorem whenever the moment power divided by
$r$ exceeds one.  The shift scale is then $n^{p/r}$, and the transformed
tail power is $(\gamma+1)/r-1$.  This is a direct change-of-variable
corollary, not a separate new universality hypothesis.

\section{Complex zeros from the real transition and positivity}
\label{pou:sec:zeros}

A useful additional consequence requires no separate deformation to a
complex saddle.  Positive polynomial coefficients provide the complex
bounds needed to promote the real limit to an analytic one.
We state the result for the general moment model of
Theorem~\ref{pou:thm:universal}.

\begin{theorem}[Complex coexistence limit and a zero array]
\label{pou:thm:zeros}
Under the hypotheses of Theorem~\ref{pou:thm:universal}, extend
$\tau_n(w)$ in \eqref{pou:eq:generalwindow} affinely to complex $w$ and define
\[
 F_n(w)=\frac{M_n(n^p\tau_n(w))}{(n^p\tau_n(w))^n}.
\]
Then $F_n(w)\to1+e^{-w}$ locally uniformly on $\mathbb C$.
For every fixed integer $j$, and all sufficiently large $n$, there is a
simple zero $c_{n,j}$ of $M_n(c)$ such that
\begin{equation}\label{pou:eq:complexzeros}
 \frac{c_{n,j}}{n^p}=\tau_{p,*}+
 \frac{(\beta+\tfrac12)\log n+\log K_{p,*}+(2j+1)\pi i}
      {d_{p,*}n}+o(n^{-1}).
\end{equation}
More precisely, in any fixed sufficiently small disk about
$w_j=(2j+1)\pi i$, the function $F_n$ has exactly one zero, counting
multiplicity, for all sufficiently large $n$; that zero tends to $w_j$.
\end{theorem}
\begin{proof}
On each fixed compact set of $w$, $\tau_n(w)\to\tau_{p,*}>0$, so $F_n$
is analytic there for all sufficiently large $n$.  Write
$\tau_n(w)=a_n+ib_n$ with $a_n>0$.
For bounded $w$ we have $b_n=O(1/n)$ and
\[
 |\tau_n(w)|-\tau_n(\Re w)
 =\frac{b_n^2}{|\tau_n(w)|+a_n}=O(n^{-2}).
\]
Therefore $|\tau_n(w)|=\tau_n(r_n)$ for a real number
$r_n=\Re w+O(1/n)$, uniformly on the compact set.
Since $M_n(c)$ has nonnegative coefficients, the triangle inequality gives
\begin{equation}\label{pou:eq:positivecomplexbound}
 |F_n(w)|\le
 \frac{M_n(n^p|\tau_n(w)|)}{(n^p|\tau_n(w)|)^n}
 =F_n(r_n).
\end{equation}
The real transition theorem bounds the right side uniformly on compact
sets.  Cauchy's integral formula now gives uniform derivative bounds on
smaller compact disks.  By equicontinuity and a diagonal subsequence
argument, every subsequence has a further subsequence converging locally
uniformly to an analytic function.  On the real axis that function must
be $1+e^{-w}$, by Theorem~\ref{pou:thm:universal}.  The identity theorem
therefore identifies every such analytic limit with $1+e^{-w}$.
This proves locally uniform convergence of the entire sequence.  The
analytic tools are the usual Cauchy estimates and uniqueness of analytic
continuation; see also \cite{dlmfcomplex}.

The limit $1+e^{-w}$ has a simple zero at each $w_j$ and no other zero in
a sufficiently small closed disk about $w_j$.  On the boundary circle its
modulus has a positive minimum.  Uniform convergence and Rouch\'e's
theorem then give exactly one zero of $F_n$ in that disk, counting
multiplicity.  Shrinking the disk proves convergence to $w_j$; count one
also proves simplicity.  The denominator defining $F_n$ is nonzero
there, so these are zeros of $M_n$.  Transforming the coordinate back to
$c$ gives \eqref{pou:eq:complexzeros}.
\end{proof}

\begin{corollary}[Central Poupard polynomial zeros]\label{pou:cor:poupardzeros}
For each fixed $j\in\Z$, $A_n(c)$ has a simple zero with
\begin{equation}\label{pou:eq:poupardzeros}
 \frac{c_{n,j}}{n^2}=\tau_*+
 \frac{\tfrac32\log n+\log K_*+(2j+1)\pi i}{d_*n}
 +o(n^{-1}).
\end{equation}
In particular, for all sufficiently large $n$, some zeros of $A_n(c)$
lie in the right half-plane.
\end{corollary}
\begin{proof}
Use $p=2$, $\beta=1$, $C=2$, $\alpha=\pi/2$ in the theorem.
The real part of $c_{n,j}/n^2$ tends to $\tau_*>0$.
\end{proof}

The theorem controls every fixed finite collection of these zeros, not
indices $j$ growing with $n$ or the global zero distribution.
Its locally uniform complex convergence is qualitative; the stated
real-axis error does not automatically become an identical complex error
bound.  Higher-order complex zero locations require an additional
quantitative argument.

\section{Inverting the central Poupard growth}\label{pou:sec:inverse}

\subsection{A canonical real interpolation}

For real $x\ge0$ define
\begin{equation}\label{pou:eq:realinterp}
 A_x=\int_0^\infty(1+y^2)^x\frac{y}{\sinh(\alpha y)}\dd y.
\end{equation}
This agrees with $a_n$ at integer $x=n$, is strictly increasing, starts
at $A_0=1$, and tends to infinity.  It therefore has a unique inverse
$x=x(Y)$ for $Y\ge1$.
The integral is in fact entire in the complex parameter $x$, since on
compact sets of $x$ its absolute value is bounded by a fixed power of
$1+y^2$ times an integrable exponential tail.  Only the positive real
inverse is used below.

\begin{lemma}\label{pou:lem:realexp}
The fixed-shift logarithmic expansion \eqref{pou:eq:logA}, with $c=1$ and
$n$ replaced by real $x\to\infty$, holds to every fixed order.
\end{lemma}
\begin{proof}
First consider one fixed odd exponential $q$ in the density.
Its saddle lies near $y=2x/(\alpha q)$.  Restrict to $y\ge\eta x$ with
$\eta>0$ smaller than the saddle location divided by $x$.
Here the generalized binomial expansion of
$(1+y^{-2})^x$ to any fixed number of terms has a controlled remainder,
because $x/y^2=O(1/x)$.  Integrating its $j$th term gives the factor
\[
 \frac{\fall xj}{j!}
 \frac{\Gamma(2x-2j+2)}{\Gamma(2x+2)}
 =\frac1{4^j j!\fall{x+\tfrac12}{j}}.
\]
The omitted lower region is exponentially small relative to the saddle
if $\eta$ is sufficiently small; a bounded $y$ region is smaller still.
The same localization bounds the integrated remainder at the desired
inverse-power order.  Thus the coefficients $C_r$ and then $\ell_r$
are unchanged for real $x$.  The other odd exponentials are exponentially
smaller.  Stirling's real expansion then gives \eqref{pou:eq:logA}.
\end{proof}

\begin{theorem}[Lambert inversion through the first two corrections]
\label{pou:thm:inverse}
Let $L=\log Y$, let $W$ be the positive real Lambert branch, and set
\begin{equation}\label{pou:eq:x0}
 w=W\left(\frac{L}{\alpha e}\right),\qquad
 x_0=\frac{L}{2w},\qquad C_0=\log\frac{8\sqrt\pi}{\alpha^2}.
\end{equation}
Define
\begin{align}
 d_0&=-\frac{\tfrac32\log x_0+C_0}{2(w+1)},\label{pou:eq:d0}\\
 d_1&=-\frac{d_0^2+\tfrac32d_0+13/24+\alpha^2/4}
                 {2x_0(w+1)}.\label{pou:eq:d1}
\end{align}
As $Y\to\infty$,
\begin{equation}\label{pou:eq:inverseexp}
 x(Y)=x_0+d_0+d_1+O\left(\frac1{x_0^2(w+1)}\right).
\end{equation}
In particular, $d_0\to-3/4$.
\end{theorem}
\begin{proof}
The defining identity for $W$ \cite{dlmfW} gives
$\log(2x_0/(\alpha e))=w$, hence
$2x_0\log(2x_0/(\alpha e))=L$.
Let $f_0(x)=2x\log(2x/(\alpha e))$.
Then $f_0'(x_0)=2(w+1)$ and $f_0''(x_0)=2/x_0$.
Lemma~\ref{pou:lem:realexp} says
\[
 \log A_x=f_0(x)+\tfrac32\log x+C_0+
             \frac{13/24+\alpha^2/4}{x}+O(x^{-2}).
\]
The displacement $d_0=O(1)$ cancels the logarithmic and constant residual
at $x_0$.  At $x_0+d_0$ the remaining term of order $1/x_0$ is
\[
 \frac{d_0^2+\tfrac32d_0+13/24+\alpha^2/4}{x_0}.
\]
Dividing its negative by $f_0'(x_0)$ gives $d_1$.
The residual after this correction is $O(x_0^{-2})$.
The derivative of $\log A_x$ is asymptotic to $2\log(2x/\alpha)$;
this follows either by differentiating the localized integral expansion
or by inserting $\log(1+y^2)$ into the saddle calculation.
The mean value theorem therefore gives the error in \eqref{pou:eq:inverseexp}.
Finally $w=\log x_0+O(1)$, so \eqref{pou:eq:d0} tends to $-3/4$.
\end{proof}

\begin{proposition}[An arbitrary-order inverse procedure]\label{pou:prop:inverseall}
For each fixed $M\ge0$, let $\xi_M$ be the sufficiently large solution of
\[
 f_0(x)+\tfrac32\log x+C_0+\sum_{r=1}^{M}h_r(1)x^{-r}=L.
\]
Then
\[
 x(Y)-\xi_M=O\left(\frac{x_0^{-M-1}}{w+1}\right).
\]
The root $\xi_M$ can be expanded to any prescribed algebraic order by
Taylor substitution about $x_0$, or computed by Newton iteration started
at $x_0+d_0$.
\end{proposition}
\begin{proof}
The truncated logarithm differs from the exact one by $O(x^{-M-1})$
and has derivative asymptotic to $2\log x$, positive for large $x$.
The same holds for the exact logarithm.  Existence, uniqueness in a
neighborhood of $x_0$, and the stated difference follow from monotonicity
and the mean value theorem.  In that neighborhood the second derivative
is $O(1/x_0)$ and the first derivative is bounded below by a positive
multiple of $w+1$.  Taylor substitution is triangular in successive
powers of $1/x_0$; Newton's quadratic error bound supplies the alternative
constructive procedure.
\end{proof}

\subsection{An exponentially small inverse displacement}

An exponential correction must again be measured relative to an exact
sector, not relative to a finitely truncated algebraic approximation.
Define
\begin{equation}\label{pou:eq:firstreal}
 A_x^{[1]}=2\int_0^\infty(1+y^2)^x y e^{-\alpha y}\dd y,
\end{equation}
and let $x_1(Y)$ be its real inverse for sufficiently large $Y$.

\begin{proposition}\label{pou:prop:inverseexponential}
As $Y\to\infty$,
\begin{equation}\label{pou:eq:inverseexponential}
 x(Y)-x_1(Y)=
 -\frac{3^{-2x_1-2}}{2\log(2x_1/\alpha)}
       \left(1+O(x_1^{-1})\right).
\end{equation}
In particular, the displacement is negative.
\end{proposition}
\begin{proof}
For each fixed odd $q$, the one-exponential integral has the leading
factor $2\Gamma(2x+2)(\alpha q)^{-2x-2}$ and correction
$1+O(1/x)$, by the proof of Lemma~\ref{pou:lem:realexp}.
The sum of the terms $q\ge5$ is $O(5^{-2x-2})$ relative to the first
integral.  To bound it directly, use
\[
 \sum_{q\ \mathrm{odd}\ge5}e^{-\alpha qy}
 =\frac{e^{-5\alpha y}}{1-e^{-2\alpha y}}
 \le e^{-5\alpha y}\left(1+\frac1{2\alpha y}\right).
\]
Both resulting integrals have the same exponential saddle factor, and
one has an additional relative factor $O(1/x)$.
Consequently
\[
 \log A_x-\log A_x^{[1]}
 =3^{-2x-2}(1+O(1/x)).
\]
The derivative of $\log A_x^{[1]}$ is
$2\log(2x/\alpha)+O(1/x)$.  Evaluating at $x_1$, and applying Taylor's
formula to the inverse, gives \eqref{pou:eq:inverseexponential}; the change
in the small correction over this displacement is of strictly smaller
order.  Positivity of all omitted sectors also gives $A_x>A_x^{[1]}$
and hence $x(Y)<x_1(Y)$.
\end{proof}

For integer enumeration the exact counting inverse is
$\max\{n:a_n\le Y\}=\lfloor x(Y)\rfloor$.
It is not valid to floor an uncertified asymptotic approximation near an
integer.  A numerical implementation needing an exact index must enclose
$x(Y)$ or compare neighboring exact terms.

\section{Reproducible exact checks and numerical diagnostics}\label{pou:sec:checks}

The accompanying program \path{code/verify.py} uses only integer arithmetic,
SymPy rational polynomial arithmetic, and mpmath numerical functions.
It checks the tangent S-fraction and the conjectured fraction through
$x^{48}$, the Poupard recurrence through row 24, the congruence through
$n=256$, and both determinant products for orders $1$ through $6$ at
$c=0,1,3$.  It also derives the correction polynomials through order eight,
checks the first saddle coefficient by symbolic differentiation, and
compares the moment integral with independent exact binomial sums.
These tests support the derivation but do not replace its all-index proofs.

All floating-point tables below use 90 decimal digits internally.  They
are reproducible diagnostics, not interval enclosures.  Exact integer
values are used for the fixed-shift and inverse tables.  The large-$n$
transition values use positive log-sum evaluation of \eqref{pou:eq:Klaw} and
\eqref{pou:eq:Tdirichlet}; they do not use numerical quadrature at a narrow
large-$n$ saddle.

\subsection{Fixed-shift accuracy}

Let $\widehat a_n^{(r)}$ be the exponential of \eqref{pou:eq:logA} truncated
after $h_r(1)n^{-r}$, and let $r=0$ mean the known leading equivalent.
Table~\ref{pou:tab:fixed} gives $|\widehat a_n^{(r)}/a_n-1|$.
The table measures a logarithmic-expansion approximation; it is not a
comparison of exact odd-pole sectors.

\begin{table}[htbp]
\centering
\caption{Relative errors for fixed-shift logarithmic truncations.}
\label{pou:tab:fixed}
\begin{tabular}{rcccc}
\toprule
$n$ & $r=0$ & $r=1$ & $r=3$ & $r=6$\\
\midrule
\tableinput{data/fixed_table.tex}
\bottomrule
\end{tabular}
\end{table}

\subsection{The transition is not already at its limit at moderate size}

The explicit window \eqref{pou:eq:window} converges relatively slowly because
its error contains $(\log n)^2/n$.  Table~\ref{pou:tab:critical} displays that
fact rather than hiding it.  Its final column is the relative error of
\[
 1+\frac2{\tau n}+
 K(\tau)n^{3/2}e^{n\Delta(\tau)}\left(1+\frac{L_1(\tau)}n\right)
\]
against the exact positive sum, evaluated at $\tau=\tau_n(w)$.
This is Theorem~\ref{pou:thm:uniform} truncated with $M=2$.

\begin{table}[htbp]
\centering
\caption{Explicit-window convergence and a two-contribution correction.}
\label{pou:tab:critical}
\begin{tabular}{rrccc}
\toprule
$n$ & $w$ & Normalized moment & $1+e^{-w}$ & Corrected relative error\\
\midrule
\tableinput{data/critical_table.tex}
\bottomrule
\end{tabular}
\end{table}

At $w=0$ the normalized moment is about $2.7155$ for $n=128$, not already
two.  The uniform corrected formula nevertheless has relative error about
$2.66\times10^{-3}$.  At $n=4096$ that error is about
$3.39\times10^{-6}$.  An exact-balance centering such as
\eqref{pou:eq:exactbalance} is preferable when the goal is a precise
finite-size coexistence parameter.

\subsection{Inverse accuracy and the exact-sector checks}

For Table~\ref{pou:tab:inverse} the input is $Y=a_n$, so the true interpolated
inverse is exactly $n$.  The reported errors are absolute errors in the
index, not relative errors in $Y$.
\begin{table}[htbp]
\centering
\caption{Absolute inverse errors at the exactly known inputs $Y=a_n$.}
\label{pou:tab:inverse}
\begin{tabular}{rccc}
\toprule
$n$ & $x_0$ & $x_0+d_0$ & $x_0+d_0+d_1$\\
\midrule
\tableinput{data/inverse_table.tex}
\bottomrule
\end{tabular}
\end{table}

The verification output additionally records, at $n=8,16,32$, the exact
integer moment divided by $F_n$, the exact first polynomial sector, the
$q=3$ sector, and the positive residual bounded by
\eqref{pou:eq:tailbound} with $Q=5$.  This is a useful way to test the
beyond-all-orders normalization without subtracting two low-order
algebraic approximations.  The symbolic and numerical results, precision,
package versions, and all displayed table data are in
\path{data/verification.json}.

\subsection{Complex roots and additional universality checks}

Table~\ref{pou:tab:zeros} follows the root associated with $w=\pi i$ in
Theorem~\ref{pou:thm:zeros}, using the affine coordinate from
\eqref{pou:eq:window}.  The coefficients supplied to reverse Horner evaluation
are exact integers.  Root solving is numerical; the normalized residuals
are below $10^{-55}$ at the working precision, but that is not an interval
certificate of root location or uniqueness.  The all-large-$n$ existence
and simplicity assertion comes from the theorem, not from these residuals.
The limiting real and imaginary parts are $0$ and $\pi$, respectively.

\begin{table}[htbp]
\centering
\caption{The zero associated with $w=\pi i$ in the coexistence coordinate.}
\label{pou:tab:zeros}
\begin{tabular}{rcc}
\toprule
$n$ & $\Re w_{n,0}$ & $\Im w_{n,0}$\\
\midrule
\tableinput{data/zeros_table.tex}
\bottomrule
\end{tabular}
\end{table}

The JSON output also tests the general transition on gamma moment models
with $(p,b)=(3/2,1/4),(2,1),(3,2)$, where $b$ is the gamma shape parameter.
These include negative and positive logarithmic displacements and an
endpoint error of order $n^{-1/2}$ when $p=3/2$.
At $n=4096$ and $w=0$, their normalized moments are approximately
$2.00822$, $2.00774$, and $2.03211$, respectively, approaching the common
limit two.  These diagnostics supplement, rather than establish,
Theorem~\ref{pou:thm:universal}.

\section{Further research questions and a formalization route}
\label{pou:sec:future}

The following questions are proposed research directions, not conclusions
of the present proofs.  They are chosen to extend a structural result or
remove a specific limitation, rather than merely ask for more terms of
a calculation whose general coefficient algorithm is already given.

\begin{question}[A direct combinatorial explanation of the new fraction]
Can one give a weight-preserving bijection from the appropriate Poupard
objects, or their marked variants, to Dyck paths with successive weights
$9,8,25,24,\ldots$?  Theorem~\ref{pou:thm:bala} proves equality through formal
contraction.  A bijection would explain the squared-odd and adjacent-even
weights at the level of objects rather than generating functions.
\end{question}

\begin{question}[Which shifts admit polynomial S-weights?]
For which real or integer shifts $c$ does the J-fraction
\eqref{pou:eq:AJ}, or a fixed Christoffel transform of its moment measure,
admit S-weights that are polynomial functions of the index on each parity
class?  The tangent and Poupard examples suggest a classification problem.
A complete answer should distinguish positivity of weights, existence of
an algebraic contraction, and a combinatorial interpretation.
\end{question}

\begin{question}[Minimal modular descriptions]
Corollary~\ref{pou:cor:modrational} gives an elementary finite-height bound.
What is the minimal rational denominator modulo $m$, and what are the
minimal preperiod and period of $A_n(c)$?  The modulus-nine collapse is
much smaller than the generic height bound.  Prime-power lifting and
cancellation of path products may produce comparably sharp formulas for
other moduli, but those formulas have not been proved here.
\end{question}

\begin{question}[A local limit theorem across coexistence]
Can the two-point weak limit and conditional Gaussian law be upgraded to
a uniform lattice local limit theorem, including the exact probability
of each fixed small index and Edgeworth corrections in the macroscopic
window?  Such a result must match an endpoint series to a lattice
Gaussian rather than approximate both components by one continuous law.
The explicit $E_m$ and $L_m$ are natural inputs but do not by themselves
establish a local theorem.
\end{question}

\begin{question}[Growing-index zero arrays and higher corrections]
Theorem~\ref{pou:thm:zeros} determines zeros corresponding to each fixed
$w_j=(2j+1)\pi i$.  How large may $j=j(n)$ grow while a uniform analogue
remains valid, and what are the next corrections to their locations?
Determine whether this local array joins a global limiting zero curve.
The positivity/normal-family proof gives qualitative convergence on
fixed compact sets, not quantitative bounds on expanding complex domains.
\end{question}

\begin{question}[The boundary $p\downarrow1$]
In Theorem~\ref{pou:thm:universal}, the critical stationary point approaches
the endpoint as $p$ tends to one.  Determine a joint $n\to\infty$,
$p\downarrow1$ scaling that connects the separated first-order transition
to its endpoint-coalescence limit.  Neither the Gaussian expansion nor
the fixed-$p$ error estimate is uniform in this joint limit.
\end{question}

\begin{question}[Slowly varying and competing tails]
Extend \eqref{pou:eq:generaltail} to a slowly varying multiplier or to several
competing exponential scales.  A multiplier depending on $n$ at the
saddle should alter the balance by its logarithm, but a theorem should
specify sufficient regularity, uniform remainder bounds, and the
possibility of more than two competing components.  The exact odd-pole
decomposition offers a concrete positive model for testing such a theory.
\end{question}

\begin{question}[Certified computation and inverse rounding]
Develop explicit constants and computable thresholds in
Theorem~\ref{pou:thm:uniform}, then combine them with interval arithmetic to
certify the coexistence parameter and integer inverse.  The sector bound
\eqref{pou:eq:tailbound} is already explicit, but the all-order Laplace
remainders in this article are asymptotic bounds with unspecified
constants.  A certified implementation must close that gap rather than
interpret high-precision diagnostics as enclosures.
\end{question}

\subsection{A realistic Lean route}

The finite algebraic layer can be separated from the analytic layer.
A first formalization target is the path recurrence for $T_n$, formal
S- and J-fraction truncations, and their coefficient stabilization.
The contraction lemma and binomial translation then give
Theorem~\ref{pou:thm:bala}.  Weighted Motzkin paths modulo nine give
Theorem~\ref{pou:thm:mod9} with a small, local divisibility argument.
These steps need no asymptotic integration.

A second layer can formalize the gamma cancellation identity, the
finite polynomial $H_n$, and the elementary odd-tail integral bound.
The positive integral representation is the bridge to that layer.
The positive-coefficient complex bound and a normal-family theorem would
then support the complex zero result.
The most substantial separate obligations are a parameter-uniform
Laplace theorem with an endpoint component, differentiation and inversion
of its asymptotics, and the conditional central limit argument.
The present Python checks do not discharge those obligations.  A future
formalization should record precisely which theorems are kernel checked,
which numerical values are interval certified, and which results remain
ordinary mathematical proofs.

\section{Conclusion}

The central Poupard sequence provides two concrete, now-proved answers to
statements displayed as conjectural in OEIS: the December 2025 continued
fraction and the uniform $3$-valuation observation.  Their proofs are
short because the right formal representation exposes the relevant
structure.

The larger contribution of the investigation is a coherent analytic
picture.  Fixed shifts possess exact positive exponential sectors with
explicit tail control.  Quadratic shifts create an endpoint--saddle
competition, a logarithmically displaced coexistence window, and a
macroscopic two-component index law.  Positivity promotes the real limit
to an analytic limit and rigorously locates nearby complex zeros.
The same mechanism persists for a
large class of exponentially decaying moment measures, while the power
in their tail determines the logarithmic displacement.  Lambert inversion
then connects the fixed-shift transseries to a usable index estimate,
without confusing algebraic truncation errors with exponentially small
sectors.

The known leading equivalent is preserved and attributed; the extensions
are accompanied by proofs, coefficient algorithms, tests, and clearly
stated limitations.  The next substantial gains would come from growing-index complex
zero arrays, joint endpoint limits, local probability laws, or fully
certified remainder estimates, rather than from finite-data extrapolation.

\appendix
\section{Coefficient generation and reproduction}\label{pou:sec:reproduce}

The formal logarithm coefficients used in the code may be generated without
symbolic logarithm expansion.  If $C_0=1$ and
$\sum C_rt^r=\exp(\sum_{r\ge1}\ell_rt^r)$, differentiating gives
\[
 \ell_r=C_r-\frac1r\sum_{j=1}^{r-1}j\ell_jC_{r-j}.
\]
Together with \eqref{pou:eq:Cformula} and \eqref{pou:eq:hgeneral}, this is an exact
rational-polynomial algorithm at every finite order.  The delivered
verification file contains $C_r$, $\ell_r$, and the logarithmic $h_r$
through $r=8$; in that file the common symbolic argument is $z=c\alpha^2$.
The notation $h_r(c)$ in the article already includes that substitution.

To reproduce the supplied full diagnostics from the package directory:
\begin{verbatim}
python -m pip install -r requirements.txt
python code/verify.py --max-n 4096 --dps 90
latexmk -pdf -interaction=nonstopmode -halt-on-error article.tex
\end{verbatim}
A shorter diagnostic run can be directed to a separate directory:
\begin{verbatim}
python code/verify.py --max-n 512 --output-dir data/quick
\end{verbatim}
The PDF can be rebuilt without rerunning Python; its generated table
fragments are included.  No external font files are bundled.
The exact tests use independent recurrences, contractions, determinant
calculations, and symbolic derivative identities where practical.
The scripts contain no theorem-prover certificates and do not label
floating-point output as a rigorous error enclosure.

\writenote{In the collection the make file is \path{code/Makefile}
and the requirements file is \path{data/requirements.txt}. Its targets assume the report root as the
working directory, so it is invoked as \texttt{make -f code/Makefile}.
The full command above rewrites the six generated files in the shipped
\texttt{data/}, because \texttt{code/verify.py} writes by default to the
\texttt{data/} directory beside \texttt{code/}; the \texttt{verify}
target does the same and also overwrites
\texttt{data/verification\_run.txt} with its standard output. Run these
commands on a copy of the report directory, or pass
\texttt{--output-dir} with a scratch directory. At intake the full
command was rerun on a copy: all six generated files were equal to the
shipped ones apart from Windows line ends, and the standard output
differed only in the line naming the output directory. That line in the
shipped \texttt{data/verification\_run.txt} names the authors' working
directory \texttt{/mnt/data/\ldots}, which is not part of the
collection.}

\section{Source and novelty audit}\label{pou:sec:audit}

The source review checked the OEIS entries available at inspection for A125054, A125053,
A000182, and A000364; Bala's continued-fraction note; the primary
Foata--Han finite-difference paper; the indicated NIST DLMF sections;
and the ProveIt repository material described in Section~\ref{pou:sec:scope}.
Targeted searches for the central Poupard asymptotics, shifted tangent
moments, and transition terminology did not identify a matching
publication for the full package of extensions.  Such a search is not
proof of publication novelty.  In particular, the proofs use standard
Laplace and continued-fraction methods, and specialized consequences may
have precedents under other terminology.

Repository code searches for A125054 and A000182 returned no matches in
the inspected search results.  Search indexing and unexamined material
prevent this from being an exhaustive nonduplication certificate.
The earlier snapshot was kept fixed even though the repository's main
branch advanced during the investigation.  Full details and source URLs
are recorded in \path{PROVENANCE.md}.

The verification boundary is exact: the continued-fraction and congruence
statements have all-index mathematical proofs; the analytic theorems have
ordinary mathematical proofs with stated asymptotic errors; finite exact
computations and floating diagnostics are supplementary.  No claim of
independent peer review, Lean verification, or interval certification is
made.

\section{Provenance of the collection edition}\label{pou:app:edition}

\writenote{This appendix was added in the write.}

\paragraph{Source.} The report is built from one manuscript, number~47
of batch~73O1 of ProveIt's incoming-reports intake: the archive
\texttt{Central\_Poupard\_OEIS.zip} (16~files, a 28-page PDF), which
arrived in commit \texttt{3a9518c52} and was placed in the collection by
commit \texttt{9df4ba51a}. Its pinned repository commit is
\texttt{ef66df0fdb8fc64cbc80756d622cab0e8e85e262} (1~October~2026);
\path{PROVENANCE.md} also records a later read of
\texttt{dc1a7242d2ab4a4496b7dfee63256ee767b21fc9}, not used for content.
The author line is kept as delivered and names no assistant; like every
manuscript of the intake, the report is AI-assisted and unrefereed.

\paragraph{What the write changed.} Every one of the 108 delivered
labels received the prefix \texttt{pou:}, and every reference was
updated with it; this appendix adds one label. The write added the
\texttt{[Write note]} paragraphs (after the abstract, in
Sections~\ref{pou:sec:scope} and~\ref{pou:sec:geometry}, in
Appendix~\ref{pou:sec:reproduce} and here), the bibliography entries
\cite{bellreport} and \cite{chartpath}, and the \verb|\writenote|
macro. The \verb|\input| and \verb|\tableinput| paths
\texttt{data/\ldots} are unchanged. No delivered sentence, statement,
proof or table was changed or removed.

\paragraph{Delivered names and shipped paths.} \texttt{article.tex},
\texttt{README.md} (rewritten as the collection's report guide) and
\path{PROVENANCE.md} are at the report root; \texttt{code/verify.py} and
the seven files of \texttt{data/} kept their paths; the
\texttt{Makefile} moved from the package root to \texttt{code/}, and
\texttt{requirements.txt} and \texttt{BUILD\_REPORT.json} from the
package root to \texttt{data/}. All but the article and the README are
byte-identical to the delivery. Not shipped: the delivered
\texttt{article.pdf} (replaced by a build of this edition) and the
checksum ledger \texttt{SHA256SUMS} (verified 15/15 at intake, then
retired). \texttt{data/BUILD\_REPORT.json} describes the delivered
28-page PDF and the delivered source, which are no longer the shipped
files.

\paragraph{Relation to repository material.} No other repository report
or formal development treats A125054 or its analytic model. The
Lambert-chart transseries package \cite{chartpath} is a methodological
precedent only, and the shared constant $s_*$ with the A277364 report
\cite{bellreport} is discussed in the notes in
Sections~\ref{pou:sec:scope} and~\ref{pou:sec:geometry}. None of the
statements of this report is formalized.

\begin{thebibliography}{99}
\small
\raggedright
\setlength{\parskip}{0pt}
\setlength{\itemsep}{0.45em}

\bibitem{oeis125054}
OEIS Foundation Inc., \emph{A125054: Central terms of triangle A125053}.
Includes the leading asymptotic attributed to V. Kotesovec (30 May 2015),
the valuation observation of F. Chapoton (2 August 2021), and the
continued-fraction conjecture of P. Bala (15 December 2025).
Inspected 1 October 2026. \url{https://oeis.org/A125054}.

\bibitem{oeis125053}
OEIS Foundation Inc., \emph{A125053: A variant of the Poupard triangle
with secant-number first column}.  Inspected 1 October 2026.
\url{https://oeis.org/A125053}.

\bibitem{oeis182}
OEIS Foundation Inc., \emph{A000182: Tangent numbers}.
Inspected 1 October 2026. \url{https://oeis.org/A000182}.

\bibitem{oeis364}
OEIS Foundation Inc., \emph{A000364: Euler (secant) numbers}.
Inspected 1 October 2026. \url{https://oeis.org/A000364}.

\bibitem{foatahan}
D. Foata and G.-N. Han, \emph{Finite Difference Calculus for Alternating
Permutations}, arXiv:1304.2483, 2013.  In particular, Theorems 1.1--1.5
and the bivariate generating function in equation (1.10).
\url{https://arxiv.org/abs/1304.2483}.

\bibitem{bala}
P. Bala, \emph{Some S-fractions related to the expansions of
$\sin(ax)/\cos(bx)$ and $\cos(ax)/\cos(bx)$}, 11 May 2017,
OEIS-hosted note. \url{https://oeis.org/A002439/a002439.pdf}.

\bibitem{dlmfcos}
NIST Digital Library of Mathematical Functions,
\emph{\S4.22: Infinite Products and Partial Fractions}.
\url{https://dlmf.nist.gov/4.22}.

\bibitem{dlmfgamma}
NIST Digital Library of Mathematical Functions,
\emph{\S5.11: Asymptotic Expansions, Gamma Function}.
\url{https://dlmf.nist.gov/5.11}.

\bibitem{dlmflaplace}
NIST Digital Library of Mathematical Functions,
\emph{\S2.3: Integrals of a Real Variable}, especially the treatment of
Laplace's method. \url{https://dlmf.nist.gov/2.3}.

\bibitem{dlmfcomplex}
NIST Digital Library of Mathematical Functions,
\emph{\S1.10: Functions of a Complex Variable}.
\url{https://dlmf.nist.gov/1.10}.

\bibitem{dlmfW}
NIST Digital Library of Mathematical Functions,
\emph{\S4.13: Lambert $W$-Function}.
\url{https://dlmf.nist.gov/4.13}.

\bibitem{repo}
V. Reshetnikov, \emph{ProveIt}, repository snapshot
\texttt{ef66df0fdb8fc64cbc80756d622cab0e8e85e262}.
\url{https://github.com/VladimirReshetnikov/ProveIt/tree/ef66df0fdb8fc64cbc80756d622cab0e8e85e262}.

\bibitem{repochart}
\emph{The Logarithmic Critical Endpoint: Convergent Lambert Charts,
All-Order Sector Laws, and Sharp Action-Cutoff Corrections},
ProveIt research report, 29 September 2026, README inspected at the
snapshot in \cite{repo}.  Methodological context only.
\url{https://github.com/VladimirReshetnikov/ProveIt/blob/ef66df0fdb8fc64cbc80756d622cab0e8e85e262/Analysis/Transseries/docs/series-and-transseries/Logarithmic_Critical_Endpoint_Lambert_Charts/README.md}.

\bibitem{chartpath}
[Added in the write.] The package of \cite{repochart} in the repository:
\path{Analysis/Transseries/docs/series-and-transseries/Logarithmic_Critical_Endpoint_Lambert_Charts/}.

\bibitem{bellreport}
[Added in the write.] ProveIt research-report collection,
\emph{Asymptotics of OEIS A277364},
\path{SetTheory/Cardinals/docs/reports/generating-functions-and-asymptotics/oeis-sequence-asymptotics/a277364-bell-asymptotics/A277364_asymptotics.tex},
equation \texttt{eq:taulambert}.

\end{thebibliography}
\end{document}
